\section{Digital compilation under slab anticoncentration}
\label{sec:slab-quantized-input}
The binary quantizer also works for singular input laws, including
surface measure on a sphere. The approximation argument here uses
output-level hyperplanes rather than translation of a density.

Assume $X\in[-1,1]^d$ and, for fixed $\kappa\ge1$ and
$0<\alpha\le1$,
\begin{equation}\label{eq:digital-slab-assumption}
 P_X\{|u\cdot X-t|\le r\}\le\kappa r^\alpha
 \qquad(\|u\|_2=1,\ t\in\mathbb R,\ r>0).
\end{equation}
No Lebesgue density is required.

\begin{lemma}[Slab-regular quantized-input compiler]
\label{lem:slab-quantized-compiler}
Let $g$ be a ReLU network of hidden width $C\ge1$, hidden depth
$h\ge1$, arbitrary finite real coefficients, and affine biases.
For $0<\epsilon<1/4$, define
\[
 A=(C+1)^{dh},\qquad L=\lceil8/\epsilon\rceil,
 \qquad F=(hC+L+1)A,
\]
\[
 J=\left\lceil\log_2\max\left\{8,
       4\sqrt d\left(\frac{8\kappa F}{\epsilon}\right)^{1/\alpha}
                          \right\}\right\rceil,
\]
and
\[
 K=4(dJ+1)(C+1)^h+\lceil\log_2(32/\epsilon)\rceil+2.
\]
For a numerical exponent $b\ge2$, and every prescribed hidden depth
$H\ge2J+b(h+1)+3$, there is a tanh network $q$ with raw affine
readout satisfying
\[
 \E|q(X)-\operatorname{clip}(g(X),[-1,1])|^2\le6\epsilon,
 \qquad \sup_{[-1,1]^d}|q|\le4.
\]
Its readout coefficient absolute sum plus absolute bias is at most
four. Its hidden width and coefficient cap satisfy
\[
 m\le c K^b,\qquad
 \log(B+2)\le\frac c\alpha
   \left[J+\log\frac{\kappa(d+1)}\epsilon+\log(K+2)\right],
\]
with numerical $c,b$ independent of the original coefficients and
the additional learner depth. In particular,
\[
 J\le\frac c\alpha
   \left[dh\log(C+1)
           +\log\frac{\kappa(d+1)(hC+1)}\epsilon\right].
\]
Thus inverse accuracy enters width and depth only logarithmically.
\end{lemma}

\begin{proof}
The common refinement of all hidden activation regions has at most
$A$ full-dimensional polyhedral cells. Indeed, each layer cuts an
existing cell by at most $C$ affine hyperplanes, producing at most
$\sum_{j=0}^d\binom Cj\le(C+1)^d$ subcells. The activation boundaries
are contained in a union of at most $hCA$ affine hyperplanes: include
each nonconstant preactivation's zero hyperplane on each preceding
cell. Identically zero or nonzero constant expressions add no boundary.
This is an upper bound on full hyperplanes, not merely on their
restricted pieces. Boundary points are contained in the same union;
degenerate identically zero gates do not create further boundaries.

On each final activation cell the raw output is affine. For every
such affine formula $\ell$, add its nonempty proper level hyperplanes
\[
 \ell(x)=-1+2j/L,\qquad j=0,\ldots,L.
\]
When $\ell$ is constant there is no level boundary to add. The total
collection $\mathcal H$ has at most $F$ hyperplanes. It depends on
the teacher, but its cardinality bound is independent of all teacher
coefficient magnitudes. On a connected set avoiding every plane of
$\mathcal H$, the activation formula stays fixed and the clipped
output varies by at most $2/L\le\epsilon/4$. If the raw output lies
beyond the extreme levels, its clipped value is constant instead.

Put $Q=2^J$, $y=(x+2\mathbf1)/4$, and
\[
 x_Q=-2\mathbf1+4Q^{-1}\lfloor Qy\rfloor.
\]
Here $y\in[1/4,3/4]^d$, and
$\|x-x_Q\|_2\le4\sqrt d/Q$.
Whenever the rounding segment intersects a plane in $\mathcal H$,
the original point lies within $4\sqrt d/Q$ of that plane.
The slab assumption and union bound therefore give an exceptional
set of probability at most
\[
 \kappa F(4\sqrt d/Q)^\alpha\le\epsilon/8.
\]
Outside that set, writing $f=\operatorname{clip}(g,[-1,1])$,
we have $|f(X)-f(X_Q)|\le\epsilon/4$. This argument uses one fixed
unshifted grid and works for the singular laws allowed by
\eqref{eq:digital-slab-assumption}.

For the binary quantizer choose
\[
 \rho=\frac14\left(\frac{\epsilon}{16\kappa d Q}\right)^{1/\alpha}.
\]
Since $\alpha\le1$, this is smaller than $1/(4Q)$.
In the range $y_j\in[1/4,3/4]$, at most $Q$ grid boundaries occur
for each coordinate. Each original-coordinate bad strip has width
condition $|x_j-t|<4\rho$. Consequently its union has probability
at most
\[
 \kappa dQ(4\rho)^\alpha=\epsilon/16.
\]
No density bound is used for this second exceptional set either.

The exact quantizer bits describe $X_Q$ as an affine function of
$dJ$ signed Boolean variables. Compose that decoding with the first
layer of $g$ and apply the coefficient-free Boolean compiler from
Lemma~\ref{lem:raw-Boolean-tanh}. It has size $S\le K^b$, depth
at most $b(h+1)$, and a signed-bit readout with variation plus offset
below four, approximating $f(X_Q)$ to $\epsilon/4$.
Use Lemma~\ref{lem:tanh-quantizer} at tolerance
\[
 \delta=\min\left\{\frac{\rho}{10^6Q(J+1)},
                   \frac1{100(S+1)},\frac\epsilon{100}\right\}.
\]
Integer-margin majority gates replaced by tanh gates restore this
accuracy at each layer. The raw readout adds at most $4\delta$,
so $|q-f(X_Q)|\le\epsilon/2$ away from the quantizer strips.
Additional hidden layers restore the signed bits before the final
readout. They do not accumulate error or increase the coefficient cap.

Outside the two exceptional sets the total prediction error is at
most $3\epsilon/4$, and hence its square is at most $\epsilon^2$.
Everywhere, $|q-f|\le5$ because the readout variation and offset
give $|q|\le4$. The mean square error is thus at most
\[
 \epsilon^2+25(\epsilon/8+\epsilon/16)\le6\epsilon.
\]
The quantizer and Boolean widths are bounded by $cK^b$.
Their logarithmic coefficient bound follows from
\[
 \log(1/\rho)
 =\log4+\alpha^{-1}\log(16\kappa dQ/\epsilon)
\]
and the explicit quantizer gains and majority-gate biases.
Finally, taking logarithms in the definition of $F$ and $J$ gives
the stated upper bound on $J$. All constants are numerical.
\end{proof}

\begin{corollary}[Uniform cubes and spheres]
\label{cor:digital-cube-sphere}
The preceding compiler applies in both of the following cases.
If $P_X$ has density at most $\kappa_0\ge1$ relative to the uniform
probability law on $[-1,1]^d$, one may take
$\alpha=1$ and $\kappa=\kappa_0\sqrt d$.
If $d\ge2$ and $P_X$ has density at most $\kappa_0$ relative to
uniform surface probability on $\mathbb S^{d-1}$, one may take
\[
 (\alpha,\kappa)=
 \begin{cases}
 (1/2,4\kappa_0),&d=2,\\
 (1,4\kappa_0\sqrt d),&d\ge3.
 \end{cases}
\]
In particular the sphere version has no ambient Lebesgue-density
assumption. The regularity constants are fixed uniformly over the
teacher and the requested accuracy.
\end{corollary}
\begin{proof}
For the uniform cube, choose a coordinate with
$|u_j|\ge1/\sqrt d$ and condition on all the others. Its uniform
coordinate lies in an interval of length at most $2r/|u_j|$,
giving slab probability at most $\sqrt d\,r$. A bounded density
relative to this reference probability multiplies the bound by
$\kappa_0$.

For the unit sphere, rotation invariance reduces the question to
the first coordinate. When $d\ge3$, its density is
\[
 c_d(1-t^2)^{(d-3)/2}\mathbf1_{\{|t|<1\}},\qquad
 c_d=\frac{\Gamma(d/2)}{\sqrt\pi\,\Gamma((d-1)/2)}\le\sqrt d.
\]
An interval of length $2r$ therefore has probability at most
$2\sqrt d\,r$. For $d=2$ the density is
$1/(\pi\sqrt{1-t^2})$ on $(-1,1)$; integrating over any interval
of length $2r$ gives at most $4\sqrt r$.
Multiplication by $\kappa_0$ and harmless enlargement give the
displayed constants.
\end{proof}
